% Trigonometrie · Winkel berechnen · Prüfungs-Fokus MSA – bank/trigonometrie/e2.jsonl (zusammenbau v0.6, Prüfungs-Fokus)
\einheitenkopf[e2]{Einheit 2 · Winkel berechnen}
\zweigzeile{ab Kl. 10 (am Gymnasium ab Kl. 9) · P10 ×9}
\setcounter{aufgabe}{0}

% Anlauf (ohne Prüfkennung)
\begin{aufgabe}{Winkel berechnen – Anlauf}
\begin{teile}
\teil Gegeben: Hypotenuse $7$ cm und Gegenkathete $4$ cm. Gesucht: der Winkel $\alpha$. Was ist gesucht? \\ \kreuz{Seite: Gleichung umstellen} \\ \kreuz{Winkel: Umkehrtaste}
\teil Gegenkathete $3$ cm, Hypotenuse $8$ cm – wie groß ist $\alpha$? $\alpha \approx$ \leerfeld $^\circ$
\teil Ankathete $5$ cm, Hypotenuse $7$ cm – wie groß ist $\alpha$? $\alpha \approx$ \leerfeld $^\circ$
\end{teile}
\end{aufgabe}

\begin{aufgabe}{Winkel berechnen – Prüfungsaufgaben}
\begin{teile}
\teil Eine Seilbahn führt von der Talstation $A$ zur Bergstation $B$. Das Seil ist $610$ m lang, die waagerechte Strecke unter dem Seil $460$ m. Berechne den Winkel zwischen Seil und Waagerechter. \hfill (P24) \\ \leerfeld $^\circ$
\teil Eine Materialseilbahn ist $245$ m lang, waagerecht überbrückt sie $190$ m. Berechne den Winkel zwischen Seil und Waagerechter. \hfill (P24) \\ \leerfeld $^\circ$
\teil Beim Handball steht die Spielerin $A$ $8$ m von der Torecke $C$ entfernt; der Punkt $B$ liegt so, dass bei $B$ ein rechter Winkel ist und $AB = 3$ m. Zeige rechnerisch, dass der Winkel $\alpha$ bei $A$ rund $68^\circ$ beträgt. \hfill (P19)
\teil Auf einem Bolzplatz ist $AB = 7$ m und $AC = 15$ m, bei $B$ liegt ein rechter Winkel. Zeige rechnerisch, dass der Winkel $\alpha$ bei $A$ rund $62^\circ$ beträgt. \hfill (P19)
\teil Im Parallelogramm $ABCD$ ist $BC = 26$ cm. $F$ liegt auf der Verlängerung von $AB$ über $B$ hinaus, $CF$ steht senkrecht auf $AF$ und $BF = 9$ cm. Zeige, dass der Winkel $\varepsilon = \angle CBF$ rund $70^\circ$ groß ist. \hfill (P26F)
\teil Im Parallelogramm $ABCD$ ist $BC = 18$ cm. $F$ liegt auf der Verlängerung von $AB$, $CF$ steht senkrecht auf $AF$ und $BF = 7$ cm. Berechne den Winkel $\varepsilon = \angle CBF$. \hfill (P26F) \\ $\varepsilon \approx$ \leerfeld $^\circ$
\teil Im Dreieck $ABC$ ohne rechten Winkel ist die Höhe von $C$ auf $AB$ $5{,}8$ m lang, $b = AC = 9{,}6$ m. Berechne den Winkel $\alpha$. \hfill (P22) \\ $\alpha \approx$ \leerfeld $^\circ$
\teil Im Dreieck $ABC$ ohne rechten Winkel ist die Höhe von $C$ auf $AB$ $3{,}9$ cm lang, $b = AC = 8{,}5$ cm. Berechne den Winkel $\alpha$. \hfill (P22) \\ $\alpha \approx$ \leerfeld $^\circ$
\teil Eine Rampe überbrückt waagerecht $240$ cm und steigt dabei $18$ cm. Berechne die Länge $x$ der Rampe und ihren Anstiegswinkel $\alpha$. \hfill (P25) \\ $x \approx$ \leerfeld[cm] , $\alpha \approx$ \leerfeld $^\circ$
\teil Eine Laderampe überbrückt waagerecht $300$ cm und steigt dabei $25$ cm. Berechne die Länge $x$ der Rampe und ihren Anstiegswinkel $\alpha$. \hfill (P25) \\ $x \approx$ \leerfeld[cm] , $\alpha \approx$ \leerfeld $^\circ$
\teil Die Punkte $A$, $B$, $C$ bilden ein Dreieck mit rechtem Winkel bei $A$. Berechne die Länge von $BC$ und die Größe des Winkels $\beta$ bei $B$ (Längen in Kästchen). \hfill (P19) \\

\begin{ksys}[xmin=-2,xmax=4,ymin=-4,ymax=3,ablesen] \punkt{-1}{-3}{A} \punkt{3}{-3}{B} \punkt{-1}{2}{C} \end{ksys}
$BC \approx$ \leerfeld , $\beta \approx$ \leerfeld $^\circ$
\teil Die Punkte $A$, $B$, $C$ bilden ein Dreieck mit rechtem Winkel bei $A$. Berechne die Länge von $BC$ und die Größe des Winkels $\beta$ bei $B$ (Längen in Kästchen). \hfill (P19) \\

\begin{ksys}[xmin=-1,xmax=4,ymin=-6,ymax=3,ablesen] \punkt{0}{2}{A} \punkt{3}{2}{B} \punkt{0}{-5}{C} \end{ksys}
$BC \approx$ \leerfeld , $\beta \approx$ \leerfeld $^\circ$
\end{teile}
\end{aufgabe}

\medskip
\uebersichtskasten{Winkel aus zwei Seiten: Bilde das Verhältnis der beiden gegebenen Seiten vom gesuchten Winkel aus (Gegenkathete oben, Hypotenuse oder Ankathete unten) und drücke die Umkehrtaste – sin⁻¹, cos⁻¹ oder tan⁻¹ macht aus dem Verhältnis den Winkel. \\ Hypotenuse 11 cm, Gegenkathete von α 7 cm: sin α = 7/11 ≈ 0,636 \quad α = sin⁻¹(7/11) ≈ 39,5° \\ Derselbe Bruch mit der Ankathete: cos β = 7/11 \quad β ≈ 50,5°  (α und β ergänzen sich zum rechten Winkel) \\ Beide Katheten 7 cm und 9 cm: tan α = 7/9 \quad α = tan⁻¹(7/9) ≈ 37,9° \\ Merkregel: Seite gesucht → Gleichung umstellen. Winkel gesucht → Verhältnis, dann Umkehrtaste; der Sinus- und der Kosinuswert sind nie größer als eins.}
